\scanpage{23}
\mspage{- IV-\handwritten{951} -}
et il est clair p\deleted{\unsure{q}}\unsure{a|q}r (18.\unsure{6}.6,(v)) que si $A$ est n\unsure{o|h}ethérien \deleted{\illegible{1}} , il en est de
même de $ {}^{\mathrm h}A$ . \deleted{Enfin , $ {}^{\mathrm h}A$ possède aussi la propriété universelle (18.6.6,(ii))
en convenant d'appeler``homomorphisme \inserted{\handwritten{semi-}}local'' de $A$ dans un anneau \inserted{\handwritten{semi-}}\uline{local} $B$ tout
homomorphisme $\handwritten{\varphi}$ tel que $\handwritten{\varphi}(\underline r)$ soit contenu dans l'idéal maximal \inserted{$\underline n$} de $B$ . En
effet , il est clair que pour un tel homomorphisme , $\handwritten{\varphi}^{-1}(\underline n)$ doit contenir
l'un des $\underline m_i$ puisqu'il contient leur intersection , donc $\handwritten{\varphi}^{-1}(\underline n)$ est égal à
l'un des $\underline m_i$ , et $\handwritten{\varphi}$ se factorise \deleted{\unsure{donc}} en $A\to A_{\underline m_i}\to B$ ; d'où l'on conclut
\deleted{\illegible{1}} que $\operatorname{Hom.loc}(A,B)=\handwritten{\prod_i}\operatorname{Hom.loc}(A_{\underline m_i},B)$ , et le fait que (18.6.6,(ii))
est encore valable lorsque $A$ est semi-local en découle aussitôt .}
% slanted
\margin{\handwritten{vider, (la forme hybride ici ne mérite pas le nom de ``\unsure{prop|propr.} universelle'') ou bien \unsure{arranger} avec \deleted{des} ``hom $1/2$ locaux'' dans \deleted{des} anneaux \unsure{henséliens}}}

\uline{Proposition} (18.6.8) .- \uline{Soient} $A$ \uline{un anneau semi-local} , $B$ \uline{une} $A$-\uline{algèbre semi-locale et entière sur} $A$ . \uline{Alors} $B\otimes_A({}^{\mathrm h}A)$ \uline{est une} $B$-\uline{algèbre semi-locale isomorphe à} $ {}^{\mathrm h}B$ .

\deleted{\illegible{7}} Comme $B\otimes_A({}^{\mathrm h}A)$ est
entier au-dessus de $ {}^{\mathrm h}A$ , chacun de ses idéaux maximaux est au-dessus d'un
des idéaux maximaux de $ {}^{\mathrm h}A$ , donc au-dessus d'un idéal maximal de $A$ , et comme
$B$ est entière sur $A$ , les idéaux premiers de $B$ au-dessus d'un idéal maximal
de $A$ sont des idéaux maximaux de $B$ , si bien que finalement un idéal maximal
de $B\otimes_A({}^{\mathrm h}A)$ a pour projections dans $\operatorname{Spec}({}^{\mathrm h}A)$ et $\operatorname{Spec}(B)$ des idéaux maximaux
; compte tenu de (18.6.6,(iii)) et de (\uline{I},3.4.9) , on conclut que l'anneau
$B\otimes_A({}^{\mathrm h}A)$ est semi-local . En outre , pour tout anneau local hensélien\inserted{$C$},
$\operatorname{Hom.loc}(B\otimes_A({}^{\mathrm h}A),C)$ est en correspondance biunivoque avec l'ensemble des
couples $(\handwritten{\varphi},\handwritten{\psi})$ d'homomorphismes locaux $\handwritten{\varphi}:B\to C$ , $\handwritten{\psi}:{}^{\mathrm h}A\to C$ tels que les composés
$A\to B\xrightarrow{\handwritten{\varphi}}C$ , $A\to{}^{\mathrm h}A\xrightarrow{\handwritten{\psi}}C$ soient égaux . Mais \deleted{\illegible{2}} en vertu
de la correspondance biunivoque entre $\operatorname{Hom.loc}(A,C)$ et $\operatorname{Hom.loc}({}^{\mathrm h}A,C)$ , on
voit que pour tout $\handwritten{\varphi\in}\operatorname{Hom.loc}(B,C)$ il existe un $\handwritten{\psi}$ et un seul ayant \deleted{\unsure{les}} \inserted{la} propriété\deleted{s}
précédente\deleted{s} , donc l'application $\operatorname{Hom}(B\otimes_A({}^{\mathrm h}A),C)\to\operatorname{Hom}(B,C)$ est bijective
, ce qui prouve la proposition en vertu de l'unicité de la solution
d'un problème universel .
% slanted
\margin{\handwritten{\uline{\unsure{cavalier|cavalé}}}}

\uline{Théorème} (18.6.9) .- \uline{Soit} $A$ \uline{un anneau semi-local} .

(i) \uline{Pour que} $ {}^{\mathrm h}A$ \uline{soit réduit} (resp. \uline{normal),il faut et il suffit que} $A$ \uline{le soi}\handwritten{t}.

(ii) \uline{Supposons} $A$ \uline{noethérien . Alors , pour tout idéal premier} $\underline p$ \uline{de} $A$ , \uline{l'an}-
% archive sheet number: 22
